% Transcription — fonds Alexandre Grothendieck, shelfmark 37, batch 3 (pages 41-60).
% Established from the facsimile archives/batches/37/batch-03.pdf (a copy of
% 37.pdf fetched through the project's relay, grothendieck-relay.onrender.com;
% PDF page k = archivists' page 40+k, no cover sheet), per the skill
% transcribe-grothendieck. The folder is 99 pages in five batches (1-20, 21-40,
% 41-60, 61-80, 81-99); this is the third, done in parallel with the others.
%
% Pass: Opus 5.5 (claude-opus-5-5), 2026-09-26 — first pass, unchecked against the pages by a human.
%
% Pages skipped: 45 (a blank cream sheet). Pages 50 and 55 are tan covers
% inscribed in pencil, top right, « SGA 5 » and « SGA 6 »: no \page marker,
% their words become section headings with a \note (hand.md §5).
% Pages noted only (not re-set): 41, 47, 49, 53, 54.
% Pages transcribed: 42-44, 46, 48, 51, 52, 56-60.
%
% Decisions, item by item, with the evidence:
% - Pages 42-44: typed leaves of a bilingual (French / English) list of the
%   terminology and notations of SGA 4, entry by entry with the exposé and
%   number (VIII 8.1 ... XIX 6.7), continuing the leaf at page 40 (batch 2),
%   which carries at its head, in his hand, « Terminologie, notations SGA 4
%   Vol. 2 / Attention, faire deux listes disjointes ! ». DECISION: transcribed
%   in full, with his marks, as an unpublished working document of the
%   re-edition. It is not the published text of SGA 4: it is a draft state
%   anterior to the split into two lists that its own heading demands (the
%   notations are still mixed with the terms, and are only ringed in pencil
%   for the second list), it is bilingual, it carries machine x-outs, typing
%   slips (« voisingge », « Anhyankar ») and ink and pencil corrections of
%   content (« universellement / universally » called in, references added
%   such as XIX 6.1, XVI 3.7, XIX 1.2, 2.1, 3.2, a cancelled triple entry for
%   the finiteness theorem). Who typed it is not established; the ink and
%   pencil interventions are taken as his (same blue ink and hand as « P =
%   ens. des nbs. premiers » and the two handwritten entries of p. 44, same
%   pencil as the heading of p. 40), which is presumable, not proven. The
%   typed ell (a chi-like glyph with a pencil « l » in the margin) is set
%   \ell; symbols the machine lacked (mu, wp, delta, the braces of A{t}) were
%   inked into blanks and are set in the formulas, noted once.
%   ALIGNMENT: batch 2 (page 40) was being written in parallel and did not
%   exist when this file was finished; its treatment of p. 40 should match.
% - Page 41: a typed leaf, upside down, carrying only the first five lines of
%   the same list in an earlier state (« topologie étale (d'un schéma) »,
%   « topos étale (d'un schéma) », « site étale d'un schéma »), abandoned;
%   nothing of his on it. Given one \note, not re-set.
% - Page 47: a typed leaf, upside down, paginated « 19 » by hand, of SGA 4,
%   Exposé IV, 9.1.13 « Questions ouvertes » (published, SGA 4 t. 1, LNM
%   269): one French \note with the marks on it (a long vertical stroke
%   cancelling the leaf, which served as scrap paper for p. 46 on its other
%   side; a bracket before « Il est clair que »). Not re-set.
% - Page 49: a typed leaf, upside down, paginated « 6 » by hand, of the
%   exposé on biextensions of SGA 7 (the text cites « VII 2.9.4 » and has
%   Corollaire 1.8, Biext^0, Biext^1, the Weil sheaf on P x P'), i.e.
%   presumably Exposé VIII of SGA 7 I (published, LNM 288); cancelled by a
%   long diagonal pencil stroke; in pencil every « 1. » of the numbering is
%   overwritten « 4. » (1.3, 1.5, 1.6, 1.7, 1.8 -> 4.x). One \note, not re-set.
% - Pages 53-54: a letter of R. Hartshorne (Harvard University letterhead),
%   dated « Aug 4, 1969 », in English, blue ink, about the IHES notes
%   « Exposé V. Systèmes projectifs J-adiques, par J.P. Jouanolou » (SGA 5):
%   which seminar they belong to, a doubt on the proof of the « Theorem of
%   Shih » A.3.1 p. 46, and remarks on pp. 38, 40, 42, 45, 47 (with an
%   argument that the artinian hypothesis of p. 38 is unnecessary), written
%   while preparing « Affine Duality and Cofiniteness ». Another person's
%   letter: one \note per page, not re-set; his annotation given in full: in
%   pencil, boxed top right of p. 53, « Faire copie pour Jouanolou », and a
%   pencil stroke in the left margin against the paragraph naming the notes.
% - Pages 46, 48 (versos of the typed leaves 47 and 49, which show through),
%   51, 52, 56-60: his own manuscript notes, transcribed in full.
%
% His pagination and numbered runs: none of his own on the manuscript
% leaves. Typed leaves paginated by hand « 19 » (p. 47) and « 6 » (p. 49);
% the Hartshorne letter's second leaf is numbered « 2 » by its writer. The
% list (pp. 42-44) is unnumbered; its « Vol. 2 » division is marked on p. 42
% before « théorème de changement de base propre » (XII).
% Dates: only « Aug 4, 1969 » (Hartshorne, p. 53).
%
% What the batch is about. The end of the SGA 4 terminology list (exposés
% VIII-XIX). Pages 46 and 48: a « Dictionnaire » between finite (sober)
% spaces X, finite ordered sets I, finite lattices sigma and finite topoi E
% (Ouv, Crib, Point, Top, I^, sigma~), the question of an « opposite topos »
% (C^ and C^v), and the characterisation of topoi of the form I^ or Top(X)
% by generating families of subobjects of the final object — related to the
% questions of SGA 4 IV 9.1 on the typed leaf he wrote on. Pages 51-52: a list
% of things to do for the re-edition of SGA 5 (Exposés I, IV-VII, XIII, XIV:
% weak Lefschetz, Q_ell-sheaves, Z_ell and Q_ell language, the relatively
% smooth case, the crystalline case and congruence formulas mod p, the note
% on Katz's theorem). Pages 56-60: the same for SGA 6 (Exp. II, VI, VII:
% Hirzebruch's multiplicative sequences, the self-intersection formula for
% Gr(X), Gr(Y), the additive structure of Gr(X~) (x) Q for a blow-up; Serre
% and Jouanolou on the title page; SGA 2 V 3.5), and a torn slip (p. 58)
% making an object E formally principal under an object-group tau(E).
%
% Notation fixed once: \ell for the machine's ell; \underline{G}_m as typed;
% on pp. 46, 48 Ouv, Crib, Top, Point, Fib as operator names, I^\wedge for his
% I^, \sigma^\sim for his sigma~.
%
% Hardest pages: 46 (a diagram page, fast ink, its upper third cancelled by
% three long strokes) and 58 (fragment, torn at the foot). Apparatus: 21
% \ill{}, 28 \uncertain{} (counted on the finished file).
%
% Readings to check first: « Fib » (p. 46, twice), « Lat fin » and « top coh »
% (p. 46), the first line of p. 48 after « suffisamment de », « trop
% particuliers » (p. 51), « une note » (p. 52), and the whole of p. 58.

\documentclass[11pt,a4paper]{article}
\input{../preamble/grothendieck}

\folder{37}
\batch{3}
\pages{41}{60}
\dating{1967-1971}
\watermark{Édition de démonstration}
\foldertitle{Rééditions SGA, EGA [SGA 4 à SGA 7, EGA I et EGA II] : tapuscrits et copies de tapuscrit annoté (s.d.), lettres (1967-1971, s.d.), notes manuscrites (s.d.)}

\begin{document}

\section*{Terminologie, notations SGA 4 (suite)}

\note{le titre est celui qu'il inscrit en tête du premier feuillet de la liste,
page 40 (lot 2) : « Terminologie, notations SGA 4 Vol. 2 / Attention, faire
\emph{deux} listes disjointes ! ». La liste, dactylographiée, donne pour chaque
entrée le terme français, le terme anglais et le renvoi (exposé et numéro de
SGA 4) ; on sépare ici les trois colonnes par des tirets. Les notations sont
encore mêlées aux termes : il les entoure au crayon, pour la seconde liste ; on
le signale par « (entouré) ». Les renvois ajoutés au crayon ou à l'encre bleue,
et les deux entrées manuscrites de la page 44, sont donnés en \add{}. Les
insertions dactylographiées dans l'interligne, qu'un trait d'encre bleue appelle
à leur place, sont données dans le texte sans marque et signalées. Le $\ell$ de
la machine est rendu $\ell$ ; les signes que la machine n'avait pas
($\mu$, $\wp$, $\delta$, les accolades de $A\{t\}$) ont été tracés à la main
dans des blancs}

\page{41}
\note{feuillet dactylographié, tête-bêche, portant seulement un premier état,
abandonné, des cinq premières lignes de la liste (page 40) : « (Sch) VII 1.1 »,
« topologie étale (d'un schéma) --- étale topology (of a scheme) », « topos
étale (d'un schéma) --- étale topos (of a scheme) », « site étale d'un schéma
--- étale site of a scheme ». Rien de sa main ; non reproduit}

\page{42}

\begin{itemize}
\item suite spectrale de descente --- spectral sequence of descent --- VIII 8.1,
XVII \ldots\ \marginal{$*$}\note{astérisque à l'encre bleue dans la marge de gauche ;
le même signe revient devant « morphisme trace » et « morphisme de changement de
base », les trois entrées renvoyant à XVII}
\item Hochschild-Serre (suite spectrale de ---) --- VIII 8.4
\item \add{$\mathbb{P}$ = ens. des nbs. premiers \quad IX 1.0.}\note{à l'encre
bleue, dans un cadre ouvert tracé entre cette ligne et la suivante, avec un trait
qui l'appelle à cette place}
\item faisceau de torsion, faisceau de $p$-torsion --- torsion sheaf, $p$-torsion
sheaf --- IX 1.1\note{« faisceau de torsion » et « torsion sheaf » sont tapés dans
l'interligne et appelés à l'encre bleue}
\item faisceau de ind-$p$-groupes, faisceau de groupes ind-finis --- sheaf of
ind-$p$-groups, sheaf of ind-finite groups --- IX 1.5
\item localement constant (faisceau ---) --- locally constant sheaf --- IX 2.0
\item constructible (faisceau d'ensembles, de groupes, de modules ---) ---
constructible \struck{\ill{}} sheaf of sets, groups or modules --- IX 2.3\note{un
mot tapé est noirci à l'encre bleue}
\item connexe par arcs (schéma ---) --- arcwise connected scheme --- IX 2.12
\item Kummer (théorie de ---) --- IX 3.2
\item (entouré) $\underline{G}_m$, $(\underline{G}_m)_X$, $\mu_n$ --- IX 3.0
\item Hilbert (théorème 90 de ---) --- IX 3.3
\item diviseur de Cartier --- Cartier divisor --- IX 3.4
\item Artin-Schreier (théorie de ---) --- IX 3.5
\item (entouré) $\wp$ --- IX 3.5
\item faisceau gratte-ciel --- skyscraper sheaf --- IX 4.1
\item morphisme trace --- trace morphism --- VIII 5.1, XVII \ldots
\item $\ell$-dimension cohomologique d'un schéma --- $\ell$-cohomological
dimension of a scheme --- X 1\note{« d'un schéma » et « of a scheme » sont tapés
dans l'interligne et appelés à l'encre bleue ; dans la marge de gauche, au
crayon, un « $\ell$ » en regard de cette ligne et de la suivante, pour le signe
de la machine}
\item (entouré) $\mathrm{cd}_{\ell}(X)$ --- X 1 ; $\mathrm{cdqc}\,X$ --- X 5.0
\item fibration élémentaire --- elementary fibration --- XI 3.1
\item bon voisinage (relativement à $S$) --- nice neighbourhood relative to $S$
--- \add{XI 3.2}\note{« XI 3.2 », tapé seul à la ligne suivante, est ramené par
un trait de crayon au bout de cette entrée}
\item théorème de comparaison --- comparison theorem --- XI 4.4
\item (entouré) $X_{c\ell}$ --- XI 4.0
\end{itemize}

\marginal{Vol. 2}\note{à l'encre bleue, dans la marge de gauche, sous un double
trait qui barre toute la largeur de la page : la coupure entre les volumes
passe devant l'exposé XII}

\begin{itemize}
\item théorème de changement de base propre --- proper base change theorem ---
XII 5.1
\item morphisme de changement de base --- base change morphism --- XII 4, XVII
\ldots
\end{itemize}

\page{43}

\begin{itemize}
\item (entouré) $\mathrm{Et}(Z)$ ($Z$, schéma) --- XII 5.9
\item Chow (lemme de ---) --- Chow lemma --- XII 7.1 \struck{Lemme de Chow}
\item Théorème de finitude --- finiteness theorem --- \add{XIV 1.1, \struck{\ill{}},
XVII \ldots ; XVI 5.1, 5.2 ; (XIX 5.1}\note{renvois ajoutés au crayon au-dessus
de la ligne, en remplacement des trois sous-entrées suivantes ; plus haut, en
haut à droite, un autre renvoi au crayon dont le chiffre romain est biffé :
\struck{\ill{}} 5.2}

\struck{a) pour morphisme propre --- proper morphism --- XIV 1.1 ; b) pour
$R^i_{!}f$ --- XVII \ldots ; c) pour un morphisme de type fini de schémas
excellents --- XIX 5.1}\note{les trois sous-entrées sont encadrées et barrées
d'un zigzag au crayon ; « proper morphism XIV 1.1 » est tapé par-dessus un
premier essai}
\item (entouré) $d(x)$ \add{$\delta(x)$, $\delta(F)$} ($x$, point d'un schéma,
$F$ faisceau) --- XIV 2\note{$\delta(x)$, $\delta(F)$ à l'encre bleue au-dessus
de $d(x)$ ; « $F$ faisceau » est tapé dans l'interligne}
\item \add{Lefschetz} (théorème de \struck{Lefschetz} affine) --- affine Lefschetz
theorem --- XIV 3.1, \add{XIX 6.1}\note{« Lefschetz » est biffé dans la ligne et
récrit au crayon dans la marge de gauche, avec un tiret au-dessus de la place
biffée : l'entrée devient « Lefschetz (théorème de --- affine) »}
\item Morphisme (universellement) $n$-acyclique ($n$, entier) --- (universally)
$n$-acyclic morphisme --- XV 1.3, 1.7\note{ici et dans les trois entrées
suivantes, « (universellement) » et « (universally) » sont tapés dans l'interligne
et appelés à l'encre bleue ; dans la marge de gauche, entourée à l'encre bleue,
une paire de parenthèses « ( ) »}
\item \struck{Morphisme universellement n-acyclique universally acyclic morphism}
\item morphisme (universellement) 1-asphérique --- (universally) 1-aspheric
morphism --- XV 1.7
\item \struck{morphisme universellement 1-asphérique universally 1-aspheric
morphism}
\item Morphisme (universellement) localement $n$-acyclique --- (universally)
locally $n$-acyclic morphism --- XV 1.11
\item morphisme (universellement) localement 1-asphérique --- (universally)
locally 1-aspheric morphism --- XV 1.11
\item (entouré) $A\{t\}$ --- XV 2.3
\item théorème de changement de base par morphisme lisse --- base change theorem
for smooth base change --- XVI 1.2
\item théorème de spécialisation pour la cohomologie --- specialization theorem
for cohomology --- XVI 2.2
\item théorème de pureté cohomologique relatif --- relative purity theorem for
cohomology --- XVI 3.6, \add{3.7}
\item théorème de pureté cohomologique absolue --- absolute purity theorem for
cohomology --- XVI 3.9, XIX \add{1.2, 2.1, 3.2}
\item couple lisse de schémas relatifs --- smooth pair of relative schemes ---
XVI 3.1\note{« de schémas relatifs » est tapé dans l'interligne et appelé à
l'encre bleue}
\item Abhyankar (lemme d'--- relatif) --- relative Anhyankar's lemma --- XVI 3.5
\item théorème de comparaison pour les schémas algébriques sur $\underline{C}$
--- comparison theorem for algebraic schemes over $\underline{C}$ --- XVI 4.1
\end{itemize}

\page{44}

\begin{itemize}
\item \struck{théorème de finitude}
\item théorème d'acyclicité locale --- local acyclicity theorem --- XV 2.1, XIX
4.1
\item couple régulier de schémas --- regular pair of schemes --- XIX 3.1
\item théorème de changement de base par morphisme régulier --- base change
theorem for regular base change morphism --- XIX 4.2
\item (entouré) $\delta(F,f,y)$ --- XIX 6.0
\item formellement séparable ($k$-algèbre ---) --- formally separable \add{$k$-}algebra
--- XIX 6.7\note{« $k$- » au crayon, entouré, au-dessus de « algebra » ; sur le
« 7 » final, un chiffre au crayon, \uncertain{9}}
\item \add{étendue algébrique --- VIII 10.2 --- algebraic spread}
\item \add{espace algébrique --- VIII 10.4 --- algebraic space}
\end{itemize}

\note{les deux dernières entrées sont écrites à la main, à l'encre bleue ; les
renvois VIII 10.2 et VIII 10.4 sont entourés, et un ovale au crayon, prolongé de
deux flèches vers la droite, les enveloppe. Le reste du feuillet est vide}

\section*{Topos finis : un dictionnaire}

\page{46}

\note{notes à l'encre sur le verso de la page 47 (SGA 4 IV 9.1.13, « Questions
ouvertes », sur les sous-topos). Le tiers supérieur est barré de trois longs
traits ; on le donne biffé, en abrégé}

\struck{$E' \rightleftarrows E''$ ($\varphi''$, $\varphi'$) ; $\lambda' :
\mathrm{id}_{E'} \to \varphi''\varphi'$, $\lambda'' : \mathrm{id}_{E''} \to
\varphi'\varphi''$.}

\struck{Esp fin $\longleftrightarrow$ ord fin ; Esp fin, topologie opposée ;
espaces topologiques finis $\longleftrightarrow$ ens. préordonnés finis, $x \mapsto
x$ muni rel. \uncertain{spéc.}, topologie $\leftrightarrow$ \ill{} ; \ill{}
$\longleftrightarrow$ ens. ordonnés finis ; ens. ordonnés finis \uncertain{ou
dual}\ill{}}\note{lecture très incertaine de tout ce bloc, écrit entre les
traits d'annulation}

Esp fin $\longleftrightarrow$ préord fin

$X = (X, T) \in$ Esp fin \uncertain{sob} $\longleftrightarrow$ ord fin $\ni (X,
\omega) = I$

$X \mapsto \mathrm{Ouv}(X)$, ens. des ouverts ; \uncertain{Lat fin} \ill{}
$\sigma$ ; $E \mapsto \mathrm{Ouv}(E)$ ; $\sigma \mapsto \sigma^{\sim}$,
\uncertain{top coh} ; \uncertain{Top fin spec}\ill{} ; $\mathrm{Top}(X, T)$ ;
$E \mapsto \mathrm{Point}(E)$ ; $(X, \omega)^{\wedge}$ ; ens. des
\uncertain{cribles}.

$x \leftarrow y$ : \uncertain{spéc.}, \uncertain{gén.}

\struck{ensemble polyèdre} schéma simplicial \add{fini} (= ens. avec ensemble
filtrant de parties.)

\struck{$x \leftarrow y$} \quad $x \to y$

\note{sur cette moitié de page, des flèches relient les objets entre eux sans
qu'on puisse toujours dire lesquelles ; on ne donne que les étiquettes lisibles,
dans l'ordre de la page}

\subsection*{Dictionnaire}

\note{titre souligné, dans un grand cadre. Quatre sommets : $X$ (espace fini), $I$
(ensemble ordonné fini), $\sigma$ (au centre) et $E$ (topos) ; on donne les
flèches une à une, avec leur étiquette}

\begin{itemize}
\item $X \to I$ : $\underline{\mathrm{Point}}(X)$ ou relation de générisation ;
$I \to X$ : topologie des cribles ;
\item $X \to \sigma$ : $\mathrm{Ouv}(X)$ ; $\sigma \to X$ :
\uncertain{$\underline{\mathrm{Fib}}(\sigma)^{\circ}$} ;
\item $I \to \sigma$ : $\mathrm{Crib}(I)$ ; $\sigma \to I$ :
\uncertain{$\underline{\mathrm{Fib}}(\sigma)^{\circ}$} ;
\item $E \to \sigma$ : $\mathrm{Ouv}(E)$ ; $\sigma \to E$ : $\sigma^{\sim}$ ;
\item $X \to E$ : $\mathrm{Top}(X)$ ; $E \to X$ : $\mathrm{Point}(E)$ ;
\item $I \to E$ : $I^{\wedge}$ ; $E \to I$ : $\underline{\mathrm{Point}}(E)$ ;
\item vers l'extérieur du cadre : de $X$, « topologie opposée » ; de $I$, « ens.
ordonné opposé » ; de $\sigma$, « \uncertain{ens. ord. opposé} » ; de $E$,
entouré : « ? $E \mapsto E^{\mathrm{opp}}$ ? ».
\end{itemize}

Hors du cadre : \struck{$x \leftrightarrow y$}

\[
I^{\wedge} \qquad I^{\vee}
\]

$C^{\wedge}$ et $C^{\vee} = C^{\circ\wedge}$ \quad \emph{topos opposé} ?
\note{« $C^{\wedge}$ et $C^{\vee} = C^{\circ\wedge}$ » encadré ; « topos opposé ? »
entouré, « opposé » souligné. L'exposant de $C^{\circ}$ est lu sous réserve}

\page{47}
\note{feuillet dactylographié, tête-bêche, paginé « 19 » à la main : SGA 4,
exposé IV, 9.1.13 « Questions ouvertes » (questions a) à d) sur le treillis des
sous-topos d'un $\mathbb{U}$-topos $E$), texte publié, non reproduit. Il sert de
brouillon à la page 46, écrite au verso. Marques : un long trait vertical au
crayon qui traverse tout le feuillet (annulation) ; un crochet devant « Il est
clair que » ; les $\cap$ et $\neq$ sont tracés à la main dans des blancs de la
machine. Rien d'autre de sa main}

\page{48}

\note{à l'encre, sur le verso de la page 49 ; suite de la page 46}

$I$ ens. ordonné.

On va interpréter $I^{\wedge}$ comme un $\mathrm{Top}(X)$.

On a $I \simeq \mathrm{Top}(X)$ ; considérons la topologie \uncertain{canonique}
$\sigma$, P. un
\uncertain{crible} de $I$ : partie de $I$ telle que pour tout $x \in U$, les
$y \leq x$ sont dans $U$. On peut munir $I$ de la topologie définie par cet ens.
de \struck{cribles} parties (\struck{anti}-\uncertain{pellicules}\ill{}, la
topologie des cribles). On trouve que \struck{\ill{}} l'\uncertain{on} \ill{}.\note{le $\sigma$ après « canonique » est écrit sous « topologie »}

\[
\mathrm{Ouv}(I) \overset{\mathrm{déf}}{=} \mathrm{Crib}(I) \simeq
\mathrm{Ouv}(I^{\wedge}) \quad \text{donc}
\]
\[
I^{\wedge} \simeq \mathrm{Ouv}(I)^{\sim} = \mathrm{Top}(I).
\]

\emph{NB} La relation $y \leq x$ signifie que $y$ est une \emph{générisation} de
$x$ ($x$ : une spécialisation).

$E$ a suffisamment de \ill{}, et est engendré par ses \uncertain{ouverts}

$\Updownarrow$

$E$ est engendré par les \struck{objets} sous-objets de l'objet final qui sont
\uncertain{connexes}, \uncertain{non vides} et projectifs.

$\Updownarrow$

(Ex. $E = C^{\wedge}$). Tous les morphismes de $C$ sont des monomorphismes.

$\Updownarrow$

(Ex. $E = \mathrm{Top}(X)$, $X$ sobre). Pour tout $x \in X$, l'ens. des
générisations de $x$ est un voisinage de $x$.

\note{les doubles flèches verticales relient, à gauche, les quatre alinéas ; la
première, simple, relie les deux premiers}

\page{49}
\note{feuillet dactylographié, tête-bêche, paginé « 6 » à la main, de l'exposé
sur les biextensions de SGA 7 (le texte renvoie à « VII 2.9.4 » ; il s'agit
vraisemblablement de l'exposé VIII de SGA 7 I, publié) : fin de la démonstration
d'une proposition sur le faisceau de Weil sur $P \times P'$, les
$\mathrm{Biext}^{i}$ ($i = 0, 1$) et la suite exacte $0 \to L \to P \to A \to
0$, puis début du « Corollaire 1.8 ». Non reproduit. Marques : un long trait
diagonal au crayon qui traverse le feuillet (annulation) ; au crayon, le « 1 »
de chaque numéro (1.3, 1.5, 1.6, 1.7, (1.7.1), 1.8) est surchargé en « 4 » ; un
trait vertical dans la marge devant le Corollaire. Il sert de brouillon à la page
48, écrite au verso}

\section*{SGA 5}

\note{inscrit au crayon en haut à droite d'une chemise de papier bistre (page 50),
qui ne porte rien d'autre ; elle ne reçoit pas de numéro de page}

\page{51}

\note{au crayon, sur un feuillet quadrillé détaché d'un bloc ; ici et dans les pages suivantes, les numéros d'exposé, soulignés par lui, sont en italique}

SGA 5 \emph{VII} \quad th. de Lefschetz faible sans \uncertain{supposer} $X$
et $Y$ lisses (seulement $X - Y$ lisse \ldots)

SGA 5 \emph{V}\uncertain{I} \quad $\mathbb{Q}_\ell$-faisceaux : signaler
définition raisonnable en termes de \emph{champ associé}\note{le chiffre romain
est repassé ; on lit VI sous réserve}

SGA 5. \struck{\ill{}} \add{\emph{VI}}\note{le premier chiffre romain est
noirci, « VI » est écrit dessous et entouré} \quad Reprendre en langage de
$\mathbb{Z}_\ell$-faisceaux et $\mathbb{Q}_\ell$-faisceaux les notions
essentielles dont il a été question dans SGA 4 (Ajouter nouveau N°).

SGA 5 \emph{IV} \quad Essayer de faire dans le cas relativement lisse toutes
les constructions essentielles

\note{trait horizontal}

\emph{VI} \quad Systèmes projectifs \add{-adiques} (dans les catégories
dérivées \ldots\ (ou dans thèse)

SGA 5 V \quad dans introduction du Hom ou de l'expo\supplied{sé}, signaler l'emploi
des coeff. \uncertain{trop} particuliers ($\mathbb{Z}_\ell$, $\mathbb{Z}/\ell^n
\mathbb{Z}$ \ldots)

\note{trait horizontal}

I §5 \quad Donner, si $i : y \to X$ inclusion du pt générique et $F_y$ Module loc.
libre sur $y$,
\[
D(i_*(F_y)) = i_*(D(F_y))
\]

\page{52}

\note{à l'encre noire, sur une feuille à en-tête de l'Institut des Hautes Études
Scientifiques}

SGA \textbf{5} \quad \emph{XIII}

Commentaires sur le cas « cristallin », notamment formules de congruence mod $p$
\note{ces trois lignes sont tenues par un crochet dans la marge de gauche}

\struck{\ill{}} \emph{IV} \quad Faire attention au cas relatif, consulter
\emph{VII} et SGA 6 \emph{X}

\emph{XIV} (Problèmes et Commentaires)

Signaler les variantes mod $p$ et variantes cristallines (cf \uncertain{une
note} à propos th. de Katz).

\page{53}
\note{lettre de R. Hartshorne, sur papier à en-tête « Harvard University,
Cambridge, Massachusetts 02138, Department of Mathematics, 2 Divinity Avenue »,
datée « Aug 4, 1969 », en anglais, à l'encre bleue, feuillet 1 sur 2, non
reproduite. Il a lu des notes de l'IHÉS intitulées « Exposé V. Systèmes
projectifs J-adiques, par J.P. Jouanolou » (SGA 5) et demande à quel séminaire
elles appartiennent et s'il en existe une édition plus récente ; il ne comprend
pas la démonstration du « Theorem of Shih », A.3.1 p. 46 (le lemme de la p. 45
ne vaudrait que pour $r > p$, cf. EGA $0_{\mathrm{III}}$ 13.5.5.2), ce qui met en
doute A.3.2 p. 47 et la Prop. 5.3.1 p. 40 — difficultés qui disparaîtraient en
travaillant dans pro-$A$ plutôt que dans la catégorie plus fine $P_{AR}$ ;
remarques sur la p. 42, ligne $-1$ (suite spectrale birégulière pour une
filtration finie) et la p. 47, ligne $-1$ (hypothèse à énoncer pour $n$ et $n+1$,
cf. EGA $0_{\mathrm{III}}$ 11.1.10). De sa main, au crayon, encadré en haut à
droite : « Faire copie pour Jouanolou » ; et un trait vertical au crayon dans la
marge de gauche, en regard du paragraphe qui nomme les notes}

\page{54}
\note{même lettre, feuillet 2 (« 2 » en tête, de la main de Hartshorne), non
reproduit : l'hypothèse artinienne de la p. 38 serait inutile si $X$ et $Z$ sont
AR-$J$-adiques noethériens, avec une démonstration par Artin-Rees dans $P_{AR}$ ;
il y est venu en préparant son article « Affine Duality and Cofiniteness ».
Signé « R. Hartshorne ». Quelques traits au crayon relient les termes des deux
suites exactes ; leur main n'est pas établie. Rien d'autre}

\section*{SGA 6}

\note{inscrit au crayon en haut à droite d'une chemise de papier bistre (page 55),
qui ne porte rien d'autre ; elle ne reçoit pas de numéro de page}

\page{56}

\note{au crayon, sur une feuille à en-tête de l'IHÉS}

SGA 6 \emph{V}\uncertain{I} \quad Donner interprétation de la relation entre
« suites multipl. de polynômes » de Hirzebruch, et séries formelles, par
l'opération $f * g$ \ldots

\note{trait horizontal}

Introduction, page de garde

Mentionner Serre et Jouanolou sur page de garde ?

\page{57}

\note{au crayon, sur une demi-feuille déchirée}

Pour Exp. II, \add{2.2.6} signaler SGA 2 V 3.5. p. 10.

\page{58}

\note{à l'encre bleue, sur une bande de papier déchirée au pied ; la dernière
ligne est en partie emportée}

Posons \quad $\bar{i}\,\bar{j} = \tau$

Pour tout $E \in \mathrm{Ob}\,C$, nous allons définir une application
\[
(*) \qquad \tau(E) \times E \xrightarrow{\ \varphi_E\ } E
\]
qui fera de $E$ un objet formellement principal sous l'objet-groupe $\tau(E)$.
\struck{Je \uncertain{prends}} \add{\ill{}}

\uncertain{Si} \ill{} \struck{\ill{}} $\ldots \to J \to E \xrightarrow{\ \varphi_E\ } B$\note{l'indice de la
flèche est lu sous réserve}

\ill{} \ill{}. \quad $\tau(E) = D$\ill{}$(J)$ \ldots

\page{59}

\note{au crayon, sur une feuille à en-tête de l'IHÉS}

SGA 6 \emph{VII} \quad marquer la formule de self-intersection pour
$\mathrm{Gr}(X)$, $\mathrm{Gr}(Y)$ (après \struck{\ill{}} \add{4.6})
\note{le numéro est surchargé ; on lit 4.6 sous réserve}

Peut-être \uncertain{bien} dans l'Exp. ?

\page{60}

\note{à l'encre, sur une feuille à en-tête de l'IHÉS}

SGA 6

Structure additive explicite de l'anneau $\mathrm{Gr}(\tilde{X}) \otimes_{\mathbb{Z}}
\mathbb{Q}$ d'un schéma éclaté \ldots\note{« Gr » est récrit sur un premier mot}

\end{document}
